% Lösungen Einheit 5 von 5 – Zusammengesetzte Figuren (zusammenbau v0.1)
\einheitenkopf{Einheit 5 von 5 · Zusammengesetzte Figuren}

\erg{26}{a) ein Rechteck und ein rechtwinkliges Dreieck \quad b) $A = 8 \cdot 3 + 3 \cdot 4 = 24 + 12 = 36$ m² \quad c) $A = 8 \cdot 3 + 8 \cdot 2{,}5 : 2 = 24 + 10 = 34$ m² \quad d) $A = 9 \cdot 7 - 4 \cdot 3 = 63 - 12 = 51$ m² \quad e) $A = 70 \cdot 45 - 15 \cdot 15 = 3\,150 - 225 = 2\,925$ cm² \quad f) $A = 90 \cdot 90 - \pi \cdot 5^2 \approx 8\,021$ cm² \quad g) $u = 2 \cdot (11 + 6) = 34$ m – die fehlende Ecke ändert den Umfang nicht \quad h) $24 \cdot 14 + \pi \cdot 7^2 \approx 336 + 153{,}9 \approx 490$ m²}
\erg{27}{a) Länge der Platte $1{,}2 : 0{,}8 = 1{,}5$ m; $1{,}5$ m $\ge 1{,}4$ m und $0{,}8$ m $\ge 0{,}75$ m – ja, sie reicht.}
\erg{28}{a) Brett $2 \cdot 0{,}25 = 0{,}5$ m²; Stücke $6 \cdot 0{,}075 = 0{,}45$ m²; Verschnitt $0{,}05 : 0{,}5 = 10\,\%$}
\erg{29}{a) Ina addiert die Teichfläche, statt sie abzuziehen. Richtig: $A = 120 - \pi \cdot 2^2 \approx 107{,}4$ m² \quad b) Ja: Die Figur bleibt dieselbe, nur die Schnittlinie ändert sich; jedes Stück wird genau einmal gezählt. \quad c) $A = 20 + 5 = 25$ m²; mit Verschnitt $25 \cdot 1{,}05 = 26{,}25$ m²; $26{,}25 \cdot 32 = 840$ €}
